
Foundation model training typically involves multiple stages---pre-training, fine-tuning, and reinforcement learning from human feedback---each requiring data curation through deduplication, perplexity filtering, and quality scoring. Current practice treats curation at each stage independently, re-optimizing thresholds without principled guidance on what transfers across stages. We introduce a transfer stability taxonomy that categorizes curation techniques by their dependence on stage-specific training objectives. Through controlled experiments across C4 pre-training data and Dolly/Alpaca instruction fine-tuning datasets using Llama-2-7B, we demonstrate categorical separation between objective-independent techniques (deduplication, perplexity filtering) and objective-dependent techniques (instruction quality filters). Objective-independent techniques transfer with 0.38\% average performance delta (95\% CI [0.30\%, 0.46\%]) and identical optimal thresholds (dedup=0.7, perplexity=500) across stages. Objective-dependent techniques show 5.04\% degradation (95\% CI [4.44\%, 5.65\%]) when transferred, with large effect size (Cohen's d=10.76) confirming categorical separation. We quantify quality-speed trade-offs for embedding-based subset selection: early-stage embeddings achieve 98\% of late-stage quality at 15\% of compute cost (6.$9\times$ speedup, 2.4\% quality penalty). These findings enable practitioners to reuse pre-training thresholds for low-level operations while focusing optimization on stage-specific strategies. All experiments were conducted at proof-of-concept level with mock evaluation; absolute performance values are predicted from literature rather than measured.
